% =============================================================
% acronyms.tex — Liste des abréviations
% Uniquement des sigles réellement employés dans les documents
% projet MENAL (HLD/LLD/méthodologie) ou introduits au chapitre 2.
% Compléter au fil de la rédaction — ne pas ajouter un sigle non
% utilisé ailleurs dans le rapport.
% =============================================================
\chapter*{Liste des acronymes}
\addcontentsline{toc}{chapter}{Liste des acronymes}
\thispagestyle{plain}

\begin{longtable}{|P{3cm}|P{12.4cm}|}
\hline
\textbf{Sigle} & \textbf{Signification} \\
\hline
\endhead
ADR & Architecture Decision Record --- registre de décisions d'architecture \\ \hline
ALB & Application Load Balancer \\ \hline
API & Application Programming Interface \\ \hline
BQ & BigQuery (service d'entrepôt de données GCP) \\ \hline
CI/CD & Continuous Integration / Continuous Deployment \\ \hline
CIS & Center for Internet Security --- référentiel de configuration sécurisée \\ \hline
CMEK & Customer-Managed Encryption Keys \\ \hline
CVE & Common Vulnerabilities and Exposures \\ \hline
DevSecOps & Development, Security and Operations \\ \hline
DNS & Domain Name System \\ \hline
EBIOS RM & Expression des besoins et identification des objectifs de sécurité --- Risk Manager, méthode d'analyse de risque orientée métier \\ \hline
GCP & Google Cloud Platform \\ \hline
GKE & Google Kubernetes Engine \\ \hline
HA & Haute disponibilité (High Availability) \\ \hline
IaC & Infrastructure as Code \\ \hline
IAM & Identity and Access Management \\ \hline
JWT & JSON Web Token \\ \hline
KMS & Key Management Service \\ \hline
NAT & Network Address Translation \\ \hline
NIST & National Institute of Standards and Technology \\ \hline
OIDC & OpenID Connect \\ \hline
ONNX & Open Neural Network Exchange \\ \hline
OWASP & Open Worldwide Application Security Project \\ \hline
PITR & Point-In-Time Recovery \\ \hline
PSA & Private Service Access \\ \hline
RBAC & Role-Based Access Control \\ \hline
RGPD & Règlement général sur la protection des données \\ \hline
SA & Service Account (compte de service) \\ \hline
SAST & Static Application Security Testing \\ \hline
SCA & Software Composition Analysis --- analyse de composition logicielle \\ \hline
SIEM & Security Information and Event Management \\ \hline
SLO & Service Level Objective \\ \hline
SLSA & Supply-chain Levels for Software Artifacts --- cadre de provenance \\ \hline
STRIDE & Spoofing, Tampering, Repudiation, Information disclosure, Denial of service, Elevation of privilege --- méthode de dérivation de menaces par flux \\ \hline
TLS & Transport Layer Security \\ \hline
VPC & Virtual Private Cloud \\ \hline
VPC-SC & VPC Service Controls \\ \hline
WAF & Web Application Firewall \\ \hline
WIF & Workload Identity Federation \\ \hline
ZTA & Zero Trust Architecture \\
\hline
\end{longtable}
% Le tableau des acronymes est hors numérotation (norme ESPRIT, C26).
% longtable incrémente le compteur "table" même sans \caption : on le rend.
\addtocounter{table}{-1}

\textit{Les familles d'identifiants employées dans le corps du mémoire --- exigences, flux,
décisions, écarts, règles de détection, tests --- sont récapitulées à la page suivante.}
